% Generated draft: requires mathematical and source review.
\documentclass[11pt]{article}
\usepackage[margin=25mm]{geometry}
\usepackage{fontspec}
\setmainfont{Latin Modern Roman}
\usepackage{amsmath,amssymb,mathtools,unicode-math}
\setmathfont{Latin Modern Math}
\usepackage{enumitem,needspace,xurl,hyperref}
\setlength{\parindent}{0pt}
\setlength{\parskip}{5pt}
\setlength{\emergencystretch}{3em}
\newcommand{\N}{\mathbb{N}}
\newcommand{\Z}{\mathbb{Z}}
\newcommand{\Q}{\mathbb{Q}}
\newcommand{\R}{\mathbb{R}}
\newcommand{\C}{\mathbb{C}}
\newcommand{\F}{\mathbb{F}}
\newcommand{\A}{\mathbb{A}}
\newcommand{\PP}{\mathbb{P}}
\newcommand{\E}{\mathbb{E}}
\newcommand{\eps}{\varepsilon}
\newcommand{\dd}{\,\mathrm{d}}
\begin{document}
\section*{UnsolvedMath ID: 7200027}
\textbf{Source problem code:} catalogue entry
\subsection{Definitions and mathematical statement}
For \(c\in\R^2\) and \(r\geq0\), let \(D(c,r)=\{x\in\R^2:|x-c|\leq r\}\) be a closed Euclidean disk. For every positive integer \(n\), determine
\[
 r(n)=\inf\{r\geq0:\ \exists c_1,\ldots,c_n\in\R^2,\quad
 D(0,1)\subseteq\bigcup_{i=1}^n D(c_i,r)\}.
\]
The problem asks for the smallest achievable common radius, as a function of \(n\), and arrangements attaining it. All covering disks have the same radius. Their centers are unrestricted, overlaps are permitted, and no requirement that a covering disk lie inside the unit disk is imposed. This uses the customary closed-disk interpretation of covering; the target is the entire disk, not only its boundary.
\subsection{Short English statement}
For each positive integer \(n\), find the least common radius of \(n\) disks whose union covers the whole unit disk.
\subsection{Sources}
\begin{itemize}
\item Local UnsolvedMath registry, permanent ID 7200027, source code catalogue entry; updated 2026-08-06; no proposer or proposal year supplied.
\item Wikipedia, \emph{List of unsolved problems in mathematics}, Geometry bullet 27, pinned revision 1366636767: \url{https://en.wikipedia.org/w/index.php?title=List_of_unsolved_problems_in_mathematics&oldid=1366636767#Geometry}. Registry access date 2026-07-29; MediaWiki extraction. Wikipedia text is CC BY-SA 4.0, with attribution and ShareAlike handling required for release. Provenance says the extracted bullet remains unpublished pending primary-source status review, although the registry published flag is true.
\item The supplied summary reports known small-\(n\) values, giving only the approximate phrase ``up to about \(n=12\) or so'', and no general resolution. It supplies no primary reference or exact verified threshold; its search claims are not checks performed for this draft.
\end{itemize}
\end{document}
